\clearpage
\section{Whole finite extensions and dyadic marking criteria}\label{app:dyadic_markings}

Evidence: direct reading of the original dyadic contract and ordinary
finite group proofs. This note does not prove a dyadic arithmetic field
construction, its simple-character existence, nontrivial stabilizer
attachment, H5, or a Lean closure. It identifies exactly what a supplied
actual host and sector must preserve, and proves the marked finite
attachment step for C=\{1\}.

\hypertarget{dyadic_markings-the-original-datum-includes-more-than-the-alternating-form}{%
\subsection{1. The original datum includes more than the alternating
form}\label{dyadic_markings-the-original-datum-includes-more-than-the-alternating-form}}

The original text
\path{closure_attempt162/inputs/original_reading.txt} states in
Definition~11.1, starting line4911, that at p=2 a restricted bridge
additionally supplies compatible quadratic refinements and an
isomorphism of the resulting Heisenberg central extensions lifting the
given symplectic isometry. Sections~13.1--13.9 define the refinements and
their Arf/Gauss signs.

The retained visual appendix is particularly explicit, at
lines6892--6947: with center \{I,-I\} fixed, squares of lifts carry the
quadratic bit. Multiplying a lift by i changes the particular finite
central-extension realization. The appendix says that allowing complex
scalar rescaling \textbf{without recording this distinction} erases data
that the bridge is required to retain.

The finite record must therefore retain at least:

\begin{itemize}
\tightlist
\item
  the labelled symplectic quotient and its pairing;
\item
  the supplied quadratic refinement q;
\item
  the specified finite central extension, its central
  embedding/character, and any marked lifts or specified extension
  isomorphism;
\item
  the resulting normalized Gauss sign, and all declared phase operators;
\item
  the complete original character values retained by the scalar pushout.
\end{itemize}

The last item alone does not retain the preceding finite-subgroup/lift
marking unless that marking is explicitly carried along. The raw text
does not require literal embeddings of the old arithmetic matrix groups
into a new type. Nor does it define a unique quadratic refinement
extracted automatically from an unmarked C× scalar extension of every
arithmetic type.

\hypertarget{dyadic_markings-general-finite-central-extension-criterion}{%
\subsection{2. General finite central-extension
criterion}\label{dyadic_markings-general-finite-central-extension-criterion}}

Let V be an elementary abelian 2-group with ordered basis
v\_1,\ldots,v\_m. Let E\_s and E\_t be central extensions of V by the
same abelian group Z. Fix the identity on Z and a quotient isomorphism
f.~Choose lifts a\_i of v\_i in E\_s and b\_i of f(v\_i) in E\_t. Write

\begin{Verbatim}[fontsize=\small,breaklines=true,breakanywhere=true]
a_i²=s_i,  b_i²=t_i,  [a_i,a_j]=c_ij, [b_i,b_j]=d_ij.
\end{Verbatim}

Every element has a unique normal form

\begin{Verbatim}[fontsize=\small,breaklines=true,breakanywhere=true]
z a_1^(epsilon_1) ... a_m^(epsilon_m), epsilon_i in {0,1},
\end{Verbatim}

and analogously on the target. The multiplication is completely
determined by the displayed squares and commutators.

\textbf{Criterion.} A center-fixing isomorphism lifting f exists
precisely when

\begin{Verbatim}[fontsize=\small,breaklines=true,breakanywhere=true]
c_ij=d_ij for every i,j,
s_i/t_i belongs to Z² for every i.                     (1)
\end{Verbatim}

Indeed, every such isomorphism must take a\_i to lambda\_i b\_i for some
lambda\_i in Z, forcing lambda\_i²=s\_i/t\_i. Conversely these equations
make the assignment a\_i -\textgreater{} lambda\_i b\_i preserve all
relations, and the normal forms prove it is bijective.

If the \textbf{particular lifts} a\_i and b\_i must correspond exactly,
the condition is s\_i=t\_i, with lambda\_i=1. If a full chosen section
is marked, rather than only generating lifts, its multiplication
coefficients must also correspond:

\begin{Verbatim}[fontsize=\small,breaklines=true,breakanywhere=true]
a(v)a(w)=c_s(v,w)a(v+w),
b(fv)b(fw)=c_t(fv,fw)b(f(v+w)),
\end{Verbatim}

and exact section preservation requires c\_s(v,w)=c\_t(fv,fw) for every
pair. Transporting the complete section under an existing group
isomorphism retains these identities automatically. Preserving just its
squares need not retain an additional independently specified section.

\hypertarget{dyadic_markings-fixed-two-element-center}{%
\subsubsection{2.1. Fixed two-element
center}\label{dyadic_markings-fixed-two-element-center}}

For Z=\{+1,-1\}, Z²=\{1\}. The lift square is independent of the lift
within its fiber and defines

\begin{Verbatim}[fontsize=\small,breaklines=true,breakanywhere=true]
a(v)²=(-1)^(q_s(v)).
\end{Verbatim}

The identity (ab)²=a²b²{[}b,a{]} gives the quadratic refinement law.
Criterion(1) becomes the exact requirement

\begin{Verbatim}[fontsize=\small,breaklines=true,breakanywhere=true]
q_t(f(v))=q_s(v) for every v.                          (2)
\end{Verbatim}

For an unspecified symplectic isomorphism between equal-dimensional
nondegenerate spaces, such an f exists exactly when their Arf invariants
agree. For a preassigned f, equality of Arf invariants alone is
insufficient: (2) must hold for that f.

\hypertarget{dyadic_markings-four-element-and-complex-centers}{%
\subsubsection{2.2. Four-element and complex
centers}\label{dyadic_markings-four-element-and-complex-centers}}

For Z=mu\_4, multiplication of a lift by i changes its square by -1. The
unmarked extension criterion preserves square classes modulo Z²=mu\_2.
In particular, squares +1 and -1 do not distinguish unmarked mu\_4
extensions. A separately marked finite mu\_2-central extension subgroup
(such as the full D8 or Q8), or prescribed lifts, still distinguishes
them.

For Z=C× every square ratio has a square root. Thus matching commutators
is sufficient for an unmarked center-fixing extension isomorphism. This
fact does not delete a marked finite subgroup from the objects being
compared.

\hypertarget{dyadic_markings-exact-d8q8-check}{%
\subsection{3. Exact D8/Q8 check}\label{dyadic_markings-exact-d8q8-check}}

In dimension two over C put

\begin{Verbatim}[fontsize=\small,breaklines=true,breakanywhere=true]
X=[[0,1],[1,0]],       Z=[[1,0],[0,-1]].
\end{Verbatim}

Then X²=Z²=I, XZ=-ZX, and {[}X,Z{]}=-I. The group generated by X and Z
has order8 and is D8. Its refinement on V=F2² is

\begin{Verbatim}[fontsize=\small,breaklines=true,breakanywhere=true]
q_D(x,y)=xy,       Arf(q_D)=0,
\end{Verbatim}

with square table (+1,+1,+1,-1) at (0,e,f,e+f), and Gauss sum +2.

The generators iX,iZ also anticommute, but each squares to -I and their
product also squares to -I. They generate Q8, with

\begin{Verbatim}[fontsize=\small,breaklines=true,breakanywhere=true]
q_Q(x,y)=xy+x+y,   Arf(q_Q)=1,
\end{Verbatim}

square table (+1,-1,-1,-1), and Gauss sum -2.

Both C× pushouts identify with the same ambient matrix group

\begin{Verbatim}[fontsize=\small,breaklines=true,breakanywhere=true]
P={z X^a Z^b : z in C×, a,b in {0,1}}.
\end{Verbatim}

Nevertheless the marked subgroups

\begin{Verbatim}[fontsize=\small,breaklines=true,breakanywhere=true]
H_D=<X,Z>,          H_Q=<iX,iZ>
\end{Verbatim}

are different. No isomorphism fixing the two-element center can map H\_D
onto H\_Q: it would preserve q and hence the Arf invariant, which
differs. There is also no center-fixing automorphism of P carrying the
marked H\_D onto H\_Q. Their unmarked pushouts are isomorphic, but the
marked objects (P,H\_D) and (P,H\_Q) are not.

Conversely a Q8 source embeds perfectly into the same P by sending its
generators to iX and iZ. This preserves every source square and its
negative Gauss sign. It does not identify those generators with the
independently marked D8 lifts X,Z. This is the distinction needed for
the arithmetic sector.

\hypertarget{dyadic_markings-marked-scalar-pushout-attachment-with-c1}{%
\subsection{4. Marked scalar-pushout attachment with
C=\{1\}}\label{dyadic_markings-marked-scalar-pushout-attachment-with-c1}}

Assume an actual type, its full character theta, and its Heisenberg
quotient W have already been constructed. Let P\_W be its actual C×
scalar pushout, acting on its actual Heisenberg representation eta\_W.
Let V\_src be a finite nondegenerate symplectic F2-space of dimension2s,
with its complete supplied dyadic marked extension data. For several
breaks use their labelled orthogonal sum, balanced central product, and
supplied product refinement.

Assume the following precise finite contract:

\begin{enumerate}
\def\labelenumi{\arabic{enumi}.}
\tightlist
\item
  The source full-character pushout P\_src has quotient V\_src and the
  same represented scalar center as P\_W.
\item
  The supplied finite central extensions and their marked lifts are
  retained as subgroups/elements of P\_src, or as already specified maps
  into P\_src. Any central character kernel already factored out must be
  recorded; a scalar pushout cannot recover invisible elements of that
  kernel.
\item
  The host has enough symplectic capacity, dim W\textgreater=2s, and the
  new attachment into this host has no independently fixed conflicting
  quadratic marking.
\item
  The declared stabilizer is C=\{1\}; the graph has one vertex.
\end{enumerate}

Then a complete marked attachment exists in the \textbf{same actual
candidate}.

\textbf{Proof.} Choose a symplectic embedding f:V\_src -\textgreater{} W
and let V be its image. The restriction P\_V of P\_W and P\_src have the
same commutator. Criterion(1) with center C× produces an isomorphism

\begin{Verbatim}[fontsize=\small,breaklines=true,breakanywhere=true]
F:P_src -> P_V
\end{Verbatim}

fixing every scalar and lifting f.~Carry every supplied finite subgroup
H\_src to the actual subgroup F(H\_src) of P\_V, and every marked lift a
to F(a). This operation does not erase square data:

\begin{Verbatim}[fontsize=\small,breaklines=true,breakanywhere=true]
F(a)²=F(a²),        [F(a),F(b)]=F([a,b]).              (3)
\end{Verbatim}

Since F fixes the scalar center, a²=(-1)\^{}q(v) implies
F(a)²=(-1)\^{}q(v). The transported refinement is therefore exactly q,
not a new choice with the same polar form. All marked section products,
specified extension maps, and finite group relations are preserved by
the same F. Consequently the Arf invariant and normalized Gauss sign are
unchanged.

Every original h in a full character domain still maps to theta\_i(h),
because F fixes the entire scalar center. Thus the dyadic marking is
retained together with the full original character values.

Write W=V perpendicular-sum U. Finite Stone--von Neumann gives

\begin{Verbatim}[fontsize=\small,breaklines=true,breakanywhere=true]
eta_W isomorphic to eta_src tensor eta_U.
\end{Verbatim}

Choose a line in eta\_U and its projector e. It lies in the actual
Heisenberg image algebra, and the marked source subgroup acts on e
eta\_W by precisely its supplied irreducible Heisenberg representation.
Its lift squares remain the operators in (3). The existing sector
recovery theorem then recovers the unchanged inducing representation
Lambda, and the same arithmetic candidate. There is no C-equivariance
equation to solve because C=\{1\}. QED.

This realizes an actual finite marked subgroup in the scalar extension
of the actual type, with actual represented operators that are complex
scalars times type operators. It does not assert that every marked
source element is the image of an unscaled element of the old or new
arithmetic matrix group. That stronger literal-embedding requirement is
absent from the present contract.

It also does not assert that the newly marked subgroup equals a
separately specified canonical native subgroup of the host. The marked
subgroup is part of the result and must be retained in the artifact.
Omitting it and recording only the unmarked P\_W would reproduce the
loss displayed by D8/Q8.

\hypertarget{dyadic_markings-if-the-host-refinement-is-independently-fixed}{%
\subsection{5. If the host refinement is independently
fixed}\label{dyadic_markings-if-the-host-refinement-is-independently-fixed}}

Suppose instead that the actual host arrives with a particular
nonsingular quadratic refinement q\_W and its actual finite mu\_2
central extension H\_W. An attachment into that marked extension
requires a quadratic isometric embedding, not merely a symplectic
embedding. There is a sharp capacity test.

\textbf{Theorem.} Let dim V\_src=2s and dim W=2R. With embeddings
otherwise free, a quadratic isometric embedding

\begin{Verbatim}[fontsize=\small,breaklines=true,breakanywhere=true]
(V_src,q_src) -> (W,q_W)
\end{Verbatim}

exists if and only if R\textgreater=s and, in the equal-capacity case
R=s,

\begin{Verbatim}[fontsize=\small,breaklines=true,breakanywhere=true]
Arf(q_src)=Arf(q_W).                                  (4)
\end{Verbatim}

If R\textgreater s, either source Arf type embeds.

\textbf{Proof.} Necessity of dimension is immediate. An
equal-dimensional embedding is an isometry, forcing (4). If
R\textgreater s, choose a nondegenerate quadratic space U of
dimension2(R-s) with

\begin{Verbatim}[fontsize=\small,breaklines=true,breakanywhere=true]
Arf(q_U)=Arf(q_W)+Arf(q_src).
\end{Verbatim}

Both Arf types exist in every positive even dimension, by original
Proposition13.7. The orthogonal sum V\_src perpendicular-sum U then has
the same dimension and Arf invariant as W. Original Standard Fact13.4
gives an isometry to W; restriction to V\_src gives the desired
embedding. QED.

After this embedding is chosen, the preimage of its image in H\_W has
exactly q\_src. The fixed-two-element-center criterion of Section~2 gives
the extension isomorphism. Generating lifts with the specified squares
can be matched. Additional fully prescribed section elements still
require their exact product relations, as described in Section~2.

Thus \textbf{one spare symplectic plane absorbs the Arf mismatch},
without changing the host's own refinement. This conclusion remains
conditional on the host's actual marked refinement/finite extension
being supplied or constructed. It does not manufacture a canonical
arithmetic quadratic datum from its dimension.

For a preassigned embedding f, the missing equation is instead

\begin{Verbatim}[fontsize=\small,breaklines=true,breakanywhere=true]
q_W(f(v))=q_src(v) for all v.                          (5)
\end{Verbatim}

Its discrepancy q\_W composed with f minus q\_src is a linear
functional. Strict capacity allows a different quadratic embedding when
it is free; it does not authorize replacing a fixed original map, fixed
depth placement, or fixed transport marking.

\hypertarget{dyadic_markings-exact-outcome-for-the-dyadic-construction}{%
\subsection{6. Exact outcome for the dyadic
construction}\label{dyadic_markings-exact-outcome-for-the-dyadic-construction}}

Under the retained one-vertex, C=\{1\}, full scalar-pushout realization
contract, an actual host with sufficient symplectic capacity does
support every supplied finite dyadic marking, provided the transported
finite subgroup, lifts, refinement, and phase are explicitly kept. The
proof is Section~4, rather than an assertion that commutators determine
the dyadic sign.

If a particular host-native quadratic extension is fixed in advance, use
Section~5. Strict capacity suffices for its quadratic embedding; equal
capacity requires the Arf equation, and a preassigned embedding requires
(5). If only the type quotient is known and a specifically canonical
host quadratic datum is requested, constructing and comparing that datum
remains an arithmetic obligation. The scalar-pushout proof must not
silently claim to have done it.

No part of this note handles nontrivial prescribed C actions. In that
case quadratic embedding, full extension action, lift markings, and
chosen honest phase operators must all be equivariant; dimension and Arf
equality alone do not settle those equations.

\hypertarget{dyadic_markings-evidence-and-author-self-review}{%
\subsection{Evidence and author
self-review}\label{dyadic_markings-evidence-and-author-self-review}}

Read directly: original Definition~11.1, Sections~13 and15, and visual
appendix pages V--22/V--23 in the retained raw file; final Block166
\path{PROOF166.md}, Section~2, and \path{LEADING_LAYER_CAPACITY166.md}, Sections~2.1
and4. The extension criterion and D8/Q8 calculation are complete finite
proofs. The quadratic embedding theorem uses the Arf classification
already explicitly supplied by the original paper. The sector conclusion
reuses the retained exact recovery theorem and does not assert new
arithmetic existence.

Self-review checked that complex lift rescaling is never treated as an
automorphism changing squares within a fixed marked extension, that the
marked subgroup remains in the output, that preassigned maps differ from
freely chosen embeddings, and that fixed native refinements require
their own actual input. No independent review of this note or new Lean
execution is claimed here.
